%% secoes/resultados/01-baixo-peso.tex — incluído por main.tex

%%---------------------------------------------------------------------%%
\subsection{Baixo peso ao nascer}\label{subsec:res-bpn}
%%---------------------------------------------------------------------%%

A Tabela~\ref{tab:bpn-dist} mostra a distribuição do baixo peso ao nascer nas oito populações. No recorte geral, a prevalência basal do desfecho foi de 7,25\%. Cada segmento tem a sua própria prevalência basal, indicada no título da respectiva tabela, e é sempre contra ela que se calcula o risco relativo (lift) dos subgrupos daquele segmento.

Chama atenção que cinco dos sete segmentos vulneráveis ficaram abaixo desse valor: indígena (6,16\%), rural (6,40\%), beneficiárias do Bolsa Família (6,54\%), quilombola (6,74\%) e municípios de até 10 mil habitantes (6,90\%). Apenas dois ficaram acima: mães com baixa escolaridade (7,90\%) e mães com menos de 15 anos (11,60\%).

\begin{table}[htbp]
\caption{Distribuição do baixo peso ao nascer por população analisada.}\label{tab:bpn-dist}%
\begin{tabularx}{\textwidth}{@{}Yrrrr@{}}
\toprule
População & Total & \makecell{Total\\(\%)} & \makecell{Nº de\\ocorrências} & \makecell{Ocorrências\\(\%)}\\
\midrule
Recorte geral & 5.242.215 & 100,00 & 380.321 & 7,25\\
Municípios até 10 mil hab. & 427.215 & 8,15 & 29.466 & 6,90\\
Rural & 1.129.845 & 21,55 & 72.315 & 6,40\\
Indígena & 46.150 & 0,88 & 2.843 & 6,16\\
Quilombola & 25.508 & 0,49 & 1.719 & 6,74\\
Mães com baixa escolaridade & 209.182 & 3,99 & 16.529 & 7,90\\
Mães menores de 15 anos & 26.652 & 0,51 & 3.091 & 11,60\\
Beneficiárias do Bolsa Família & 1.221.833 & 23,31 & 79.942 & 6,54\\
\bottomrule
\end{tabularx}
%\fontetabela
\end{table}


A Tabela~\ref{tab:bpn} apresenta os cinco subgrupos mais de maior WRAcc para o baixo peso ao nascer considerando os oito recortes deste estudo.
Como pode ser observado, alguns itens se repetem como marcadores de risco de baixo peso, com intensidades distintas entre os recortes. O número de consultas de pré-natal é o item mais recorrente, presente em 23 dos 40 subgrupos e em todos os recortes. Isolada, a faixa de 1 a 3 consultas varia de 24\% acima da basal na população indígena a 87\% entre mães com menos de 15 anos, passando por 66\% no recorte geral, e é o subgrupo de maior risco da tabela no recorte geral, entre mães com baixa escolaridade e entre beneficiárias do Bolsa Família. As regras com a faixa de 4 a 6 consultas que não incluem a primeira gestação ficam entre 20\% e 46\% acima da basal, com exceção do subgrupo 8 dos municípios de até 10 mil habitantes (57\%).
 

\begin{table}[p] % [p] coloca a tabela grande em uma página dedicada, evitando quebrar o fluxo do texto
\centering
\caption{Subgrupos baixo peso ao nascer (prevalência basal: 7,25\%).}
\label{tab:bpn}

% Ajustes para caber em uma única página com legibilidade clínica/epidemiológica
\scriptsize % Fonte padrão para tabelas densas em periódicos de saúde
\setlength{\tabcolsep}{4pt} % Reduz margem lateral interna das células (padrão é 6pt)
\renewcommand{\arraystretch}{1.08} % Compacta levemente a altura das linhas sem esmagar o texto

\begin{tabularx}{\textwidth}{@{}c Y r r r r r@{}}
\toprule
\textbf{\#} & \textbf{Subgrupo} & \textbf{Total} & \textbf{\makecell[r]{Total\\(\%)}} & \textbf{\makecell[r]{Nº de\\ocorr.}} & \textbf{\makecell[r]{Ocorr.\\(\%)}} & \textbf{Lift} \\
\midrule
\secaotabela{Recorte geral}
1 & nascimento hospitalar; consultas: 4--6 & 1.609.468 & 30,7 & 144.933 & 38,1 & 1,24 \\
2 & destino do lixo: coletado; sexo: F & 1.941.005 & 37,0 & 162.198 & 42,6 & 1,15 \\
3 & consultas: 1--3 & 427.185 & 8,1 & 51.522 & 13,5 & 1,66 \\
4 & escoamento sanitário: rede coletora; consultas: 4--6 & 606.962 & 11,6 & 64.446 & 16,9 & 1,46 \\
5 & água canalizada: sim; sexo: F; consultas: 4--6 & 604.665 & 11,5 & 62.542 & 16,4 & 1,43 \\

\secaotabela{Municípios de até 10 mil habitantes}
6 & consultas: 4--6 & 121.922 & 28,5 & 11.137 & 37,8 & 1,32 \\
7 & sexo: F & 208.607 & 48,8 & 16.217 & 55,0 & 1,13 \\
8 & consultas: 4--6; início do pré-natal: 1--3; Bolsa Família: não & 44.649 & 10,5 & 4.849 & 16,5 & 1,57 \\
9 & primeira gestação: sim & 119.156 & 27,9 & 9.814 & 33,3 & 1,19 \\
10 & primeira gestação: sim; consultas: 4--6 & 31.344 & 7,3 & 3.754 & 12,7 & 1,74 \\

\secaotabela{População rural}
11 & nascimento hospitalar; consultas: 4--6 & 374.735 & 33,2 & 28.818 & 39,9 & 1,20 \\
12 & primeira gestação: sim & 276.899 & 24,5 & 22.404 & 31,0 & 1,26 \\
13 & sexo: F & 552.138 & 48,9 & 39.698 & 54,9 & 1,12 \\
14 & primeira gestação: sim; consultas: 4--6 & 89.166 & 7,9 & 9.620 & 13,3 & 1,69 \\
15 & consultas: 1--3 & 112.897 & 10,0 & 10.701 & 14,8 & 1,48 \\

\secaotabela{População indígena}
16 & primeira gestação: sim & 8.489 & 18,4 & 734 & 25,8 & 1,40 \\
17 & sexo: F & 22.546 & 48,9 & 1.584 & 55,7 & 1,14 \\
18 & despesa com aluguel: 0; idade materna menor que 20 anos & 9.549 & 20,7 & 769 & 27,0 & 1,31 \\
19 & consultas: 1--3 & 9.933 & 21,5 & 758 & 26,7 & 1,24 \\
20 & sexo: F; mãe solteira & 8.922 & 19,3 & 694 & 24,4 & 1,26 \\

\secaotabela{População quilombola}
21 & primeira gestação: sim & 5.924 & 23,2 & 536 & 31,2 & 1,34 \\
22 & nascimento hospitalar; consultas: 1--3 & 2.902 & 11,4 & 305 & 17,7 & 1,56 \\
23 & primeira gestação: sim; consultas: 4--6 & 2.188 & 8,6 & 243 & 14,1 & 1,65 \\
24 & idade materna menor que 20 anos & 5.299 & 20,8 & 449 & 26,1 & 1,26 \\
25 & despesa com transporte: 0; Bolsa Família: não & 10.951 & 42,9 & 818 & 47,6 & 1,11 \\

\secaotabela{Mãe com baixa escolaridade}
26 & consultas: 1--3 & 28.030 & 13,4 & 3.259 & 19,7 & 1,47 \\
27 & primeira gestação: sim & 47.950 & 22,9 & 4.732 & 28,6 & 1,25 \\
28 & sexo: F & 102.821 & 49,2 & 9.055 & 54,8 & 1,11 \\
29 & consultas: 4--6; material do domicílio: tijolo/alvenaria & 49.452 & 23,6 & 4.737 & 28,7 & 1,21 \\
30 & consultas: 4--6; sexo: F & 37.176 & 17,8 & 3.708 & 22,4 & 1,26 \\

\secaotabela{Mãe com menos de 15 anos}
31 & consultas: 1--3 & 3.083 & 11,6 & 670 & 21,7 & 1,87 \\
32 & parto normal & 15.516 & 58,2 & 2.051 & 66,4 & 1,14 \\
33 & parto normal; consultas: 4--6 & 6.091 & 22,9 & 887 & 28,7 & 1,26 \\
34 & mãe solteira; parto normal; consultas: 1--3; idade materna menor que 20 anos; primeira gestação: sim & 1.178 & 4,4 & 310 & 10,0 & 2,27 \\
35 & despesa com gás de 45 ou mais; parto normal & 5.288 & 19,8 & 776 & 25,1 & 1,27 \\

\secaotabela{Beneficiário do Programa Bolsa Família}
36 & consultas: 4--6 & 388.589 & 31,8 & 31.209 & 39,0 & 1,23 \\
37 & sexo: F & 599.793 & 49,1 & 43.913 & 54,9 & 1,12 \\
38 & consultas: 1--3 & 92.550 & 7,6 & 10.359 & 13,0 & 1,71 \\
39 & consultas: 4--6; parto cesárea & 153.193 & 12,5 & 13.982 & 17,5 & 1,39 \\
40 & consultas: 4--6; sexo: F; destino do lixo: coletado & 126.825 & 10,4 & 12.141 & 15,2 & 1,46 \\
\bottomrule
\end{tabularx}
%\fontetabela
\end{table}


%geral
%No recorte geral (Tabela~\ref{tab:bpn}), o número de consultas de pré-natal está presente em quatro dos cinco subgrupos, em três deles na faixa de 4 a 6 consultas (subgrupos 1, 4 e 5) e em um na faixa de 1 a 3 (subgrupo 3). O sexo feminino aparece em dois subgrupos (2 e 5), e os demais itens ocorrem uma vez cada.

%O subgrupo 3 (Tabela~\ref{tab:bpn}), definido apenas por 1 a 3 consultas, apresenta o maior risco do desfecho da tabela, 66\% acima da basal do segmento, e a menor participação nas ocorrências (13,5\%). No extremo oposto, o subgrupo 2 registra o menor risco do desfecho (15\% acima da basal) e a maior participação nas ocorrências (42,6\%). Os subgrupos 4 e 5, ambos com 4 a 6 consultas somadas a itens de infraestrutura, ocupam posição intermediária, com risco 46\% e 43\% acima da basal.

%Como as regras são conjunções de atributos, esses riscos relativos não isolam o efeito de nenhum item individual. A comparação entre os subgrupos 1, 4 e 5, que compartilham a faixa de 4 a 6 consultas mas diferem nos demais itens, mostra risco 24\%, 46\% e 43\% acima da basal, o que indica que parte da variação decorre dos itens acompanhantes.

%Três dos cinco subgrupos incluem itens que descrevem condições de infraestrutura domiciliar presentes: coleta de lixo (subgrupo 2), escoamento por rede coletora (subgrupo 4) e água canalizada (subgrupo 5).

%municipios até 10k
%Já no recorte dos municípios de até 10 mil habitantes, o número de consultas aparece em três, sempre na faixa de 4 a 6, e a faixa de 1 a 3 não figura entre os cinco primeiros. A primeira gestação aparece em dois subgrupos (9 e 10).

%O subgrupo 10 apresenta o maior risco do desfecho da tabela, 74\% acima da basal, e a menor participação nas ocorrências (12,7\%); o subgrupo 7, definido apenas pelo sexo feminino, tem o menor risco do desfecho (13\% acima da basal) e a maior participação nas ocorrências (55,0\%).

%Os subgrupos 9 e 10 permitem uma comparação com apenas um item de diferença: o subgrupo 9 é definido somente por primeira gestação e está 19\% acima da basal; o subgrupo 10 acrescenta a faixa de 4 a 6 consultas e chega a 74\% acima da basal. A participação nas ocorrências, porém, cai de 33,3\% para 12,7\%.

%O subgrupo 8 combina a faixa de 4 a 6 consultas com início do pré-natal na categoria 1--3 e ausência de Bolsa Família, com risco 57\% acima da basal.

%Rural
%Na População rural o número de consultas aparece em três, em dois na faixa de 4 a 6 (subgrupos 11 e 14) e em um na faixa de 1 a 3 (subgrupo 15). A primeira gestação aparece em dois subgrupos (12 e 14).

%O maior risco do desfecho é o do subgrupo 14 (69\% acima da basal), que também tem a menor participação nas ocorrências (13,3\%); o menor é o do subgrupo 13, definido apenas pelo sexo feminino (12\% acima da basal), com a maior participação nas ocorrências (54,9\%).

%Assim como no segmento anterior, os subgrupos 12 e 14 diferem por um único item: a primeira gestação isolada está 26\% acima da basal, e a adição da faixa de 4 a 6 consultas eleva o risco para 69\% acima da basal, reduzindo a participação nas ocorrências de 31,0\% para 13,3\%.

%Neste segmento, a regra que contém a faixa de 1 a 3 consultas (subgrupo 15, 48\% acima da basal) apresenta risco inferior ao da regra com 4 a 6 consultas do subgrupo 14 (69\% acima da basal). Como as duas regras diferem também nos demais itens, a comparação entre as faixas não é direta.

%População indígena
%Na população indígena o sexo feminino aparece em dois subgrupos (17 e 20) e os demais itens ocorrem uma vez cada. O número de consultas figura em apenas um subgrupo, na faixa de 1 a 3 (subgrupo 19).

%O subgrupo 16, definido apenas por primeira gestação, apresenta o maior risco do desfecho da tabela (40\% acima da basal). Entre os subgrupos apresentados para todos os segmentos, esse é o maior risco observado para a regra composta unicamente por primeira gestação, acima dos 34\% do segmento quilombola e dos 26\% do rural.

%O menor risco do desfecho é o do subgrupo 17 (14\% acima da basal), que também tem a maior participação nas ocorrências (55,7\%). A menor participação nas ocorrências é a do subgrupo 20 (24,4\%), o maior valor mínimo entre os oito conjuntos analisados.

%O risco associado à faixa de 1 a 3 consultas neste segmento (24\% acima da basal) é o menor observado para essa faixa entre todos os subgrupos apresentados.

%Quilombora
%Na População quilombola a primeira gestação aparece em dois subgrupos (21 e 23) e o número de consultas também em dois (subgrupos 22 e 23), em faixas distintas. Os demais itens ocorrem uma vez cada.

%O subgrupo 23 apresenta o maior risco do desfecho (65\% acima da basal) e a menor participação nas ocorrências (14,1\%) da tabela. O subgrupo 25 tem o menor risco do desfecho (11\% acima da basal) e a maior participação nas ocorrências (47,6\%).

%Os subgrupos 21 e 23 diferem por um item: a primeira gestação isolada está 34\% acima da basal, e a adição da faixa de 4 a 6 consultas eleva o risco para 65\% acima da basal, com queda da participação nas ocorrências de 31,2\% para 14,1\%.

%Neste segmento, a regra com 4 a 6 consultas (subgrupo 23, 65\% acima da basal) apresenta risco superior ao da regra com 1 a 3 consultas (subgrupo 22, 56\% acima da basal). Como os itens acompanhantes diferem entre as duas regras, a comparação entre as faixas não é direta.

%Mães com baixa escolaridade
%Nas mães com baixa escolaridade, o número de consultas aparece em três, em um na faixa de 1 a 3 (subgrupo 26) e em dois na faixa de 4 a 6 (subgrupos 29 e 30). O sexo feminino aparece em dois subgrupos (28 e 30).

%O subgrupo 26, definido apenas pela faixa de 1 a 3 consultas, tem o maior risco do desfecho (47\% acima da basal) e a menor participação nas ocorrências (19,7\%). O subgrupo 28 tem o menor risco do desfecho (11\% acima da basal) e a maior participação nas ocorrências (54,8\%).

%As regras com 4 a 6 consultas apresentam risco 21\% e 26\% acima da basal, inferior ao da regra com 1 a 3 consultas (47\%).

%Nas mães menores de 15 anos o parto normal aparece em quatro dos cinco subgrupos (32, 33, 34 e 35) e o número de consultas em três (31, 33 e 34). Este é também o segmento com o maior baixo peso ao nascer entre os oito analisados (Tabela~\ref{tab:bpn-dist}).

%O subgrupo 34 apresenta o maior risco entre todos os subgrupos do baixo peso ao nascer (127\% acima da basal, isto é, 2,27 vezes a basal), com a menor participação nas ocorrências da tabela (10,0\%). 
%Trata-se de uma conjunção de cinco itens: mãe solteira, parto normal, 1 a 3 consultas, idade inferior a 20 anos e primeira gestação. 
%O item relativo à idade é redundante dentro de um segmento composto por mães com menos de 15 anos.

%O subgrupo 31, definido apenas pela faixa de 1 a 3 consultas, está 87\% acima da basal, o maior valor observado para essa faixa isolada entre todos os segmentos. O subgrupo 32, definido apenas por parto normal, tem o menor risco do desfecho (14\% acima da basal) e a maior participação nas ocorrências de toda a análise (66,4\%).

%As regras com 4 a 6 consultas apresentam risco 26\% acima da basal, inferior ao das regras com 1 a 3 consultas (87\% e 127\%).

%Por fim, nos beneficiárias do Bolsa Família o número de consultas aparece em quatro dos cinco subgrupos, em três na faixa de 4 a 6 (36, 39 e 40) e em um na faixa de 1 a 3 (subgrupo 38), e o sexo feminino aparece em dois (37 e 40). Este é também o segmento com o maior número de nascimentos (1.221.833).

%O subgrupo 38 tem o maior risco do desfecho (71\% acima da basal) e a menor participação nas ocorrências (13,0\%); o subgrupo 37 tem o menor risco do desfecho (12\% acima da basal) e a maior participação nas ocorrências (54,9\%). As regras com 4 a 6 consultas apresentam risco 23\%, 39\% e 46\% acima da basal, todos inferiores ao da regra com 1 a 3 consultas.

%O conjunto de itens que compõem os cinco subgrupos coincide parcialmente com o observado no recorte geral, onde também figuram o número de consultas, o sexo feminino e a coleta de lixo.
 
%Reunidos os oito recortes, alguns itens se repetem como marcadores de risco de baixo peso, com intensidades distintas entre os recortes. O número de consultas de pré-natal é o item mais recorrente, presente em 23 dos 40 subgrupos e em todos os recortes. Isolada, a faixa de 1 a 3 consultas varia de 24\% acima da basal na população indígena a 87\% entre mães com menos de 15 anos, passando por 66\% no recorte geral, e é o subgrupo de maior risco da tabela no recorte geral, entre mães com baixa escolaridade e entre beneficiárias do Bolsa Família. As regras com a faixa de 4 a 6 consultas que não incluem a primeira gestação ficam entre 20\% e 46\% acima da basal, com exceção do subgrupo 8 dos municípios de até 10 mil habitantes (57\%).
 
%subgrupos. Isolada, seu risco cresce de 19\% acima da basal nos municípios de até 10 mil habitantes para 25\% entre mães com baixa escolaridade, 26\% na população rural, 34\% na quilombola e 40\% na indígena, o maior valor dessa regra entre todos os recortes. Combinada à faixa de 4 a 6 consultas, ela forma o subgrupo de maior risco de três recortes: 74\% nos municípios pequenos, 69\% no rural e 65\% no quilombola.
 
O sexo feminino aparece em seis dos oito recortes, ausente apenas na população quilombola e entre mães menores de 15 anos. Isolado, define um subgrupo em cinco recortes (municípios de até 10 mil habitantes, rural, indígena, mães com baixa escolaridade e beneficiárias do Bolsa Família) e, em todos eles, é o subgrupo de menor risco da tabela (entre 11\% e 14\% acima da basal) e de maior cobertura (cerca de 49\% dos nascimentos e 55\% das ocorrências). No recorte geral, combinado à coleta de lixo, ocupa a mesma posição (15\% acima da basal e 37,0\% dos nascimentos). Quando acompanhado de outros itens, o risco é maior: 26\% com mãe solteira na população indígena, 26\% com 4 a 6 consultas entre mães com baixa escolaridade, e 43\% e 46\% com 4 a 6 consultas e itens de infraestrutura no recorte geral e entre beneficiárias do Bolsa Família.
 
Os extremos estão em dois recortes. Entre as mães com menos de 15 anos, segmento de maior prevalência basal (11,60\%), está o subgrupo de maior risco do desfecho (mãe solteira, parto normal, 1 a 3 consultas e primeira gestação, 127\% acima da basal, ou 2,27 vezes), além do maior risco isolado da faixa de 1 a 3 consultas (87\%) e da maior participação de uma única regra nas ocorrências (parto normal, com 66,4\% das ocorrências). No extremo oposto, a população indígena tem o menor risco máximo entre os oito recortes (40\% acima da basal) e o menor valor da faixa de 1 a 3 consultas (24\%); é também o recorte em que a menor participação nas ocorrências é a mais alta entre os oito (24,4\%).

\FloatBarrier
